\chapter{Mathematical Expressions}
For many values (boundary conditions, load values) it is possible not
only to prescribe a constant value but also a complex mathematical
expression. CFS++ utilizes for this a customized version of the {\em
  muParser} library (\url{http://muparser.sourceforge.net/}).

\noindent
Therefore one can easily
\begin{itemize}
\item prescribe analytical given, space-dependent, boundary- and load
  conditions, like e.g. a parabolic flow profile or a linear varying load
\item use time and frequency dependent variables to incorporate time-/
  frequency-dependent values
\item generate an analytical time-function, like $sin(\omega t
  + \phi)$ without actually loading it from an external timedata file
  with discrete values
\end{itemize}

\noindent
Currently the following xml-elements allow the prescription of mathematical
expressions:
\begin{itemize}
\item \verb#<dirichletInhom># for the \verb#value# and \verb#phase# attribute
\item \verb#<load># for the \verb#value# and \verb#phase# attribute
\item \verb#<regionLoad># in the mechanic PDE for the attribute \verb#value#
\end{itemize}

\section{Built-in expressions and operators}
{\em muParser} already comes with an extensive set of pre-defined
expressions and and functions (see following two tables).

Note that it is currently not possible to define own variables within
the parser, as otherwise any occurrence of a not explicitly defined
variable would create a new one, initialized with zero.
\begin{table}[ht]
 
  \begin{longtable}{c|c|c}
    \hline\noalign{\smallskip}
    name & \# of args. & description \\
    \noalign{\smallskip}\hline\noalign{\smallskip} 
    \verb#sin#     & 1 & sine function \\
    \verb#cos#     & 1 & cosine function \\
    \verb#tan#     & 1 & tangens function \\
    \verb#asin#    & 1 & arcus sine function \\
    \verb#acos#    & 1 & arcus cosine function \\
    \verb#atan#    & 1 & arcus tangens function \\
    \verb#sinh#    & 1 & hyperbolic sine function \\
    \verb#cosh#    & 1 & hyperbolic cosine \\
    \verb#tanh#    & 1 & hyperbolic tangens function \\
    \verb#asinh#   & 1 & hyperbolic arcus sine function \\
    \verb#acosh#   & 1 & hyperbolic arcus tangens function \\
    \verb#atanh#   & 1 & hyperbolic arcur tangens function \\
    \verb#log2#    & 1 & logarithm to the base 2 \\
    \verb#log10#   & 1 & logarithm to the base 10 \\
    \verb#log#     & 1 & logarithm to the base 10 \\
    \verb#ln#      & 1 & logarithm to base e (2.71828...) \\
    \verb#exp#     & 1 & e raised to the power of x \\
    \verb#sqrt#    & 1 & square root of a value \\
    \verb#sign#    & 1 & sign function -1 if $x<0$; 1 if $x>0$ \\
    \verb#rint#    & 1 & round to nearest integer \\
    \verb#abs#     & 1 & absolute value \\
    \verb#if#      & 3 & if ... then ... else ... \\
    \verb#min#     & var. & min of all arguments \\
    \verb#max#     & var. & max of all arguments \\
    \verb#sum#     & var. & sum of all arguments \\
    \verb#avg#     & var. & mean value of all arguments
  \end{longtable}
 \caption{Built-in functions}  
\end{table}

\begin{table}[ht]
  \begin{longtable}{c|c|c}
    \hline\noalign{\smallskip}
    operator & \ meaning & priority \\
    \noalign{\smallskip}\hline\noalign{\smallskip}
    \verb#=#                  & assignment      & -1 \\
    \verb#and#                & logical and      & 1 \\
    \verb#or#                 & logical or       & 1 \\
    \verb#xor#                & logical xor      & 1 \\
    \verb#<=# or \verb#le# *  & less or equal    & 2 \\
    \verb#>=# or \verb#ge# *  & greater or equal & 2 \\
    \verb#!=#                 & not equal        & 2 \\
    \verb#==#                 & equal            & 2 \\
    \verb#># or \verb#gt# *   & greater than     & 2 \\
    \verb#<# ot \verb#lt# *   & less than        & 2 \\
    \verb#+#                  & addition         & 3 \\
    \verb#-#                  & subtraction      & 3 \\
    \verb#*#                  & multiplication   & 4 \\
    \verb#/#                  & division         & 4 \\
    \verb#^#                  & raise x to the power of y & 5 \\
    \verb#_pi#                & $\pi = 3.1415926..$ & \\
    \verb#_e#                 & Eulerian number $e = 2.7182818..$ 
  \end{longtable}
  \vspace*{0.5cm}
  {\footnotesize * Due to the restriction in the xml language definition no '$<$' or
    '$>$' signs are allowed in tokens. Therefore the latter form is the
    only one possible in .xml-files. If scripting is used, both forms are
    possible. }
 \caption{Built-in operators and constants}  
\end{table}
\clearpage
\section{CFS custom variables}
In addition to the built-in variables, CFS++ registers additionally
some variables, which can be used in the mathematical expression (see
the following table).
\begin{table}[ht]
  \begin{longtable}{c|p{0.4\textwidth}|p{0.4\textwidth}}
    \hline\noalign{\smallskip}
     name & value & occurrence \\
    \noalign{\smallskip}\hline\noalign{\smallskip}
    \verb#t#     & Current time value       & transient analysis \\
    \verb#f#     & Current frequency value  & harmonic analysis \\
    \verb$x y z$ & Current global coordinate value & boundary
    conditions \newline load \\
    \verb#rad phi ax# & Current local coordinate value & regionLoad
    with cylindrical coordinate system
  \end{longtable}
  \caption{CFS specific variables}
 \end{table}

\section{Examples}
If one wants to prescribe a linear space dependent load, the following can be
written:
\begin{verbatim}
<bcsAndLoads>
  <load name="loadNode" value="10 * x" phase="cos(y)"/>
<bcsAndLoads>
\end{verbatim}
Here \verb#x# is used to denote the x-coordinate of each node
contained in the node list.
\newline

\noindent
An even easier way to define a time dependent variable than using the
\verb#dynamics# keyword is by using directly the
\verb#sin#-function. To define a variable with a sine of 1 kHz one can
simply write:
\begin{verbatim}
<bcsAndLoads>
  <dirichletInhom name="bcNode" value="10 * sin(2 * _pi * 1e-3 * t)"/>
<bcsAndLoads>
\end{verbatim}
The variable \verb#t# reflects the current time value in a transient
simulation.
Of couse one can also combine space- and time dependency:
\begin{verbatim}
<bcsAndLoads>
  <dirichletInhom name="bcNode" value="10 * x * sin(2 * _pi * 1e-3 * t)"/>
<bcsAndLoads>
\end{verbatim}
